\documentclass[12pt]{article}
\usepackage{amsmath,amssymb}

\begin{document}
\section*{{2019 AIME II Problems/Problem 14}}
Source: AoPS Wiki

\subsection*{Solution}
We will do casework on $n \pmod{5}$ and see which values work. Case 1: $n \equiv 1 \pmod{5}$ $n$ would have to be greater than $91$ because if it wasn't then we could just add denominations of $5$ until we got to $91.$ However, if $n$ is more than $91,$ then all the numbers from $91$ to $n$ which are not multiples of $5$ are also unrepresentable, which violates the problem's condition of $91$ being the largest unrepresentable number. So there are no values for this case. Case 2: $n \equiv 2 \pmod{5}$ $n + 1 \equiv 3 \pmod{5}$ is the smallest formable number 3 mod 5, so $n + 1 - 5 = n - 4$ is the largest unformable number 3 mod 5. $2n + 2 \equiv 1 \pmod{5}$ is the smallest formable number 1 mod 5, so $2n + 2 - 5 = 2n - 3$ is the largest unformable number 1 mod 5. Finally, $2n \equiv 4 \pmod{5}$ is the smallest formable number 4 mod 5, so $2n - 5$ is the largest unformable number 4 mod 5. Therefore, the largest unformable number overall is $2n - 3 \equiv 1 \pmod{5}.$ This satisfies the problem's condition! $2n - 3 = 91 \implies n = 47.$ Case 3: $n \equiv 3 \pmod{5}$ $n + 1 \equiv 4 \pmod{5},$ which is the smallest representable number 4 mod 5, so the largest unrepresentable number 4 mod 5 is $n + 1 - 5 = n -4.$ Similarly, $2n \equiv 1 \pmod{5}$ is the smallest formable number 1 mod 5 and $3n + 3 \equiv 2 \pmod{5}$ is the smallest formable number 2 mod 5 (isn't it $2n + 1$ ? ~Towerrocks), so the largest number 2 mod 5 that is unrepresentable is $3n + 3 - 5 = 3n - 2,$ but the largest number 1 mod 5 that is unrepresentable is $2n -5.$ $3n - 2 > 2n - 5,$ so there would be a larger unrepresentable number than $91$ if we were to set $2n - 5$ equal to it, which violates the problem's condition. And we can't set $3n - 2$ equal to $91$ because $91 \not\equiv 2 \pmod{5},$ so there are no valid values for this case. Case 4: $n \equiv 4 \pmod{5}$ $n + 1$ would be a multiple of $5,$ so we would be left with denominations of $5$ and $n.$ Then it would just be a straightforward application of the Chicken Mcnugget Theorem: $5n - n - 5 = 91 \implies n = 24,$ which is indeed $4 \pmod{5}.$ So this value works. Case 5: $n$ is a multiple of 5 In this case we’re left with denominations of $5$ and $n + 1,$ so by the Chicken Mcnugget Theorem we get $5(n + 1) - n - 6 = 91 \implies n = 23,$ which is not a multiple of $5.$ So there are no valid values here. The two valid values we found were $47$ and $24,$ so the answer is $47 + 24 = \boxed{71}.$ ~ grogg007

\end{document}
